\documentclass[10pt]{article}
\usepackage{scimisc-cv}
\usepackage[hidelinks]{hyperref}

\title{Curriculum Vitae -- Wonsub Cho}
\author{Wonsub Cho}
\date{}

\hypersetup{
  pdftitle={Curriculum Vitae -- Wonsub Cho},
  pdfauthor={Wonsub Cho},
  pdfsubject={Theoretical particle physics, dark matter, and early-Universe cosmology}
}
\urlstyle{same}

\cvname{Wonsub Cho}
\cvtagline{Theoretical Particle Physicist}
\cvpersonalinfo{\href{mailto:wikicho1674@gmail.com}{wikicho1674@gmail.com} \cvinfosep \href{https://wikicho.github.io}{wikicho.github.io} \cvinfosep \href{https://github.com/wikicho}{github.com/wikicho} \cvinfosep Suwon, Republic of Korea}

\begin{document}

\makecvtitle

\section{Research Profile}

Ph.D. candidate in theoretical particle physics studying dark-matter phenomenology and early-Universe cosmology.
My research develops particle-physics scenarios for dark matter and dark radiation and studies their cosmological evolution, relic abundance, and observational implications.
I have five first-author peer-reviewed publications spanning sterile-neutrino and boosted dark matter, dark-matter--neutrino interactions, Pauli-blocking stabilization, and cosmological tensions.

\section{Education}

\cvsubsection{Sungkyunkwan University (SKKU)}[Ph.D. Candidate in Physics][Theoretical Astroparticle Physics and Cosmology Group][Mar 2023--Present]
\cvdetail{Advisor}{Prof. Ki-Young Choi}

\cvsubsection{Sungkyunkwan University (SKKU)}[M.S. in Physics][Theoretical Astroparticle Physics and Cosmology Group][Aug 2020--Aug 2022]
\cvdetail{Thesis}{\emph{Effect of Boosted Dark Matter by Cosmic Particles}}

\cvsubsection{Sungkyunkwan University (SKKU)}[B.S. in Physics][Department of Physics][Mar 2014--Feb 2020]
\cvdetail{Thesis}{\emph{Searching for Boosted Dark Matter by High-Energy Neutrinos}}

\section{Research Experience}

\cvsubsection{CERN}[Graduate Research Visitor][Korea--CERN Cooperation Program, Theoretical Physics][2026]
\cvdetail{Selection}{Selected for a six-month graduate research placement for collaborative research with CERN theorists, supported by the National Research Foundation of Korea and the Korean Physical Society.}

\section{Selected Research Contributions}

\begin{itemize}
  \item Developed an inelastic-dark-matter and dark-radiation framework based on a $U(1)_D$ gauge symmetry to address the $H_0$, $S_8$, and Lyman-$\alpha$ cosmological tensions.
  \item Studied dark-matter stabilization through Pauli blocking in a degenerate fermion background using quantum-field-theory evolution equations and comparison with the Boltzmann treatment.
  \item Investigated energy transfer from the cosmic neutrino background to galactic dark-matter halos, demonstrating a mechanism for core formation and deriving constraints on dark-matter--neutrino scattering.
\end{itemize}

\section{Publications}

\begin{enumerate}
  \cvpublication{Stable Dark Matter from Pauli Blocking in the Degenerate Fermion Background with Quantum Field Theory}{\textbf{Wonsub Cho}, Ki-Young Choi, Junghoon Joh, and Osamu Seto}{\emph{Physics of the Dark Universe} \textbf{49} (2025) 101978.}{https://doi.org/10.1016/j.dark.2025.101978}{https://arxiv.org/abs/2407.08229}
  \cvpublication{Reconciling Cosmological Tensions with Inelastic Dark Matter and Dark Radiation in a $U(1)_D$ Framework}{\textbf{Wonsub Cho}, Ki-Young Choi, and Satyabrata Mahapatra}{\emph{Journal of Cosmology and Astroparticle Physics} \textbf{09} (2024) 065.}{https://doi.org/10.1088/1475-7516/2024/09/065}{https://arxiv.org/abs/2408.03004}
  \cvpublication{Cored Dark Matter Halos in the Cosmic Neutrino Background}{\textbf{Wonsub Cho}, Ki-Young Choi, and Hee Jung Kim}{\emph{Journal of Cosmology and Astroparticle Physics} \textbf{07} (2023) 013.}{https://doi.org/10.1088/1475-7516/2023/07/013}{https://arxiv.org/abs/2204.01431}
  \cvpublication{Sterile Neutrino Dark Matter with Dipole Interaction}{\textbf{Wonsub Cho}, Ki-Young Choi, and Osamu Seto}{\emph{Physical Review D} \textbf{105} (2022) 015016.}{https://doi.org/10.1103/PhysRevD.105.015016}{https://arxiv.org/abs/2108.07569}
  \cvpublication{Searching for Boosted Dark Matter Mediated by a New Gauge Boson}{\textbf{Wonsub Cho}, Ki-Young Choi, and Seong Moon Yoo}{\emph{Physical Review D} \textbf{102} (2020) 095010.}{https://doi.org/10.1103/PhysRevD.102.095010}{https://arxiv.org/abs/2007.04555}
\end{enumerate}

\newpage

\section{Selected Talks and Posters}

\cvdateditem{Aug 2025}{Poster}{``Reconciling Cosmological Tensions with Inelastic Dark Matter and Dark Radiation,'' 29th International Summer Institute on Phenomenology of Elementary Particle Physics and Cosmology (SI 2025).}
\cvdateditem{Apr 2025}{Talk}{``Reconciling Cosmological Tensions with Inelastic Dark Matter and Dark Radiation in a $U(1)_D$ Framework,'' 7th CUBES Workshop, Gurye, Republic of Korea.}
\cvdateditem{2025}{Talk}{``Reconciling Cosmological Tensions with Inelastic Dark Matter and Dark Radiation,'' Spring Meeting of the Korean Physical Society.}
\cvdateditem{2024}{Talk}{``Inelastic Dark Matter, the $H_0$ Tension, and Detection Prospects,'' Spring Meeting of the Korean Physical Society.}
\cvdateditem{Jun 2023}{Talk}{``Cored Dark Matter Halos in the Cosmic Neutrino Background,'' XVI International Conference on Interconnections between Particle Physics and Cosmology (PPC 2023), Daejeon, Republic of Korea.}
\cvdateditem{2023}{Talk}{``Cored Dark Matter Halos in the Cosmic Neutrino Background,'' Spring Meeting of the Korean Physical Society.}
\cvdateditem{2021}{Talk}{``Sterile Neutrino Dark Matter with Dipole Interaction,'' Fall Meeting of the Korean Physical Society.}
\cvdateditem{2020}{Talk}{``Searching for Boosted Dark Matter Mediated by a New Gauge Boson,'' Fall Meeting of the Korean Physical Society.}
\cvdateditem{2019}{Poster}{``Searching for Boosted Dark Matter by Cosmic Rays,'' Fall Meeting of the Korean Physical Society.}

\section{Teaching Experience}

\cvdateditem{2025--26}{Quantum Field Theory I and II}{Teaching Assistant, SKKU.}
\cvdateditem{Fall 2024}{Electromagnetism}{Teaching Assistant, SKKU.}
\cvdateditem{Fall 2023}{Relativity}{Teaching Assistant, SKKU.}
\cvdateditem{2023}{Particle Physics Summer Camp}{Teaching Assistant, KIAS.}
\cvdateditem{2020--22}{Mathematical Physics}{Teaching Assistant, SKKU.}
\cvdateditem{2020, 2023}{General Physics}{Teaching Assistant, SKKU.}

\section{Awards and Fellowships}

\cvdateditem{2020, 23, 25}{Outstanding Presentation Award}{Korean Physical Society.}
\cvdateditem{2021--22}{BK21 Innovation Research Fellowship}{}
\cvdateditem{2020--24}{Graduate Merit Scholarship}{Department of Physics, SKKU.}
\cvdateditem{2019}{Excellence Award for Undergraduate Poster Presentation}{Korean Physical Society.}
\cvdateditem{2018}{National Science and Technology Scholarship}{Korea Student Aid Foundation.}

\section{Research and Technical Skills}

\cvskill{Scientific computing}{Python, Mathematica, FeynCalc, C, and C++}
\cvskill{Research methods}{Relic-abundance calculations, cosmological evolution, and particle-physics model building}
\cvskill{Tools}{Git and \LaTeX}
\cvskill{Languages}{Korean and English}

\section{Industry Experience}

\cvsubsection{ASML Korea}[Technical Support Engineer][Metrology and Sensor Competency Engineer][Jul 2022--Oct 2022]

\end{document}
